% ============================================================================
% WORKBOOK (MODIFIED) — Section 8.4: Comparing Populations with Measures
% CCSS 7.SP.B.4 (AK/OK/TX: extended measures emphasis)
% Practice-focused reinforcement with IQR focus, multiple measures
% ============================================================================
\section{Comparing Populations with Measures of Center and Variability}
\topicTitle{Comparing Populations with Measures of Center and Variability}

\begin{quickReview}{Measures of Center and Variability}
\begin{itemize}[leftmargin=*]
    \item \textbf{Mean}: Add all values, divide by the count.
    \item \textbf{Median}: The middle value when data is ordered.
    \item \textbf{Range}: Highest $-$ lowest. Shows total spread.
    \item \textbf{IQR} (Interquartile Range): $Q_3 - Q_1$. Spread of the middle $50\%$.
\end{itemize}

Use \textbf{multiple measures} to compare groups. IQR is often more useful than range because it ignores outliers.

\medskip
\textbf{Example:} Crew~A: mean $140$, range $35$, IQR $15$. Crew~B: mean $145$, range $110$, IQR $30$. Crew~B catches more on average, but Crew~A is far more consistent.
\end{quickReview}

% ── Warm-Up ──────────────────────────────────────────────────────────────────
\begin{practiceBox}[funGreen]{\faSun~Warm-Up}

\practiceHeader[funGreen]{Find the measures.}

\prob Data: $12, 15, 18, 20, 25$. Find the mean. \answerBlank[3cm]
\answerExplain{$18$}{$(12 + 15 + 18 + 20 + 25) \div 5 = 90 \div 5 = 18$.}

\prob Data: $12, 15, 18, 20, 25$. Find the median. \answerBlank[3cm]
\answerExplain{$18$}{The middle value (3rd of 5) is $18$.}

\prob Data: $12, 15, 18, 20, 25$. Find the range. \answerBlank[3cm]
\answerExplain{$13$}{$25 - 12 = 13$.}

\prob Data: $8, 10, 14, 18, 22, 26, 30$. Find $Q_1$ and $Q_3$. \answerBlank[3cm]
\answerExplain{$Q_1 = 10$, $Q_3 = 26$}{Lower half: $8, 10, 14$. $Q_1 = 10$. Upper half: $22, 26, 30$. $Q_3 = 26$.}

\prob Data: $8, 10, 14, 18, 22, 26, 30$. Find the IQR. \answerBlank[3cm]
\answerExplain{$16$}{$\text{IQR} = Q_3 - Q_1 = 26 - 10 = 16$.}

\end{practiceBox}

% ── Practice A: Calculating and Comparing ────────────────────────────────────
\begin{practiceBox}{Calculating and Comparing}

\practiceHeader{Find the measures for each group and compare.}

\prob Group~A: $55, 60, 65, 70, 75$. Group~B: $40, 55, 65, 75, 90$. Find the mean and range of each group.
\answerExplain{A: mean $= 65$, range $= 20$; B: mean $= 65$, range $= 50$}{A: $(55+60+65+70+75) \div 5 = 65$, range $= 75-55 = 20$. B: $(40+55+65+75+90) \div 5 = 65$, range $= 90-40 = 50$. Same mean, but B is much more spread out.}

\prob Group~X: $10, 12, 14, 15, 18, 20, 22$. Group~Y: $5, 10, 14, 15, 16, 20, 30$. Find the IQR of each.
\answerExplain{X: IQR $= 8$; Y: IQR $= 10$}{X: $Q_1 = 12$, $Q_3 = 20$. IQR $= 8$. Y: $Q_1 = 10$, $Q_3 = 20$. IQR $= 10$. Group X is slightly more consistent.}

\prob Data set: $5, 8, 10, 12, 15, 18, 50$. Find range and IQR. Which better represents the spread?
\answerExplain{Range $= 45$; IQR $= 10$; IQR is better}{Range: $50 - 5 = 45$. $Q_1 = 8$, $Q_3 = 18$. IQR $= 10$. The outlier $50$ inflates the range. IQR shows the typical spread more accurately.}

\prob Two students track daily steps. Student~X: mean $8{,}000$, IQR $1{,}500$. Student~Y: mean $8{,}200$, IQR $4{,}000$. Compare.
\answerExplain{Y averages slightly more but is much less consistent}{Y's IQR is nearly triple X's, meaning Y's daily steps vary far more. X walks a more consistent amount each day.}

\end{practiceBox}

% ── Practice B: When to Use Which Measure ────────────────────────────────────
\begin{practiceBox}[funTeal]{When to Use Which Measure}

\practiceHeader[funTeal]{Decide which measure is most appropriate.}

\prob A data set has an extreme outlier. Should you use mean or median to describe the center? Explain.
\answerExplain{Median}{Outliers pull the mean toward extreme values. Median is resistant to outliers and better represents the typical value.}

\prob Two runners have identical medians but different ranges. What does this tell you?
\answerExplain{Their typical performances are similar, but one is more variable}{Same median means the middle performance is the same. Different range means one runner has more extreme best/worst times.}

\prob Why is IQR often better than range for comparing two groups?
\answerExplain{IQR ignores outliers}{Range uses only the two extreme values and can be misleading with outliers. IQR measures the middle $50\%$ and gives a stable comparison.}

\prob A teacher says ``Class~A did better than Class~B.'' What two measures would you check before agreeing?
\answerExplain{Center (mean or median) and spread (IQR or range)}{Comparing only the center is incomplete. You also need to see how spread out the scores are to determine if the difference is meaningful.}

\end{practiceBox}

% ── Practice C: True or False ────────────────────────────────────────────────
\begin{practiceBox}[funPurple]{True or False?}

\trueOrFalse{If two groups have the same mean, their data distributions must be the same.}
\answerExplain{False}{$\{0, 50, 100\}$ and $\{48, 50, 52\}$ both have mean $50$ but very different spreads.}

\trueOrFalse{IQR measures the spread of the middle $50\%$ of the data.}
\answerExplain{True}{IQR $= Q_3 - Q_1$ captures the range of the middle half of the data.}

\trueOrFalse{A larger IQR means the group is more consistent.}
\answerExplain{False}{A larger IQR means the data is more spread out. Smaller IQR means more consistency.}

\trueOrFalse{Range is always larger than or equal to IQR.}
\answerExplain{True}{Range spans all the data ($\max - \min$). IQR spans only the middle half ($Q_3 - Q_1$). The full span is always $\geq$ the middle span.}

\end{practiceBox}

% ── Word Problems ────────────────────────────────────────────────────────────
\begin{practiceBox}[funBlue]{\faGlobe~Word Problems}

\wordProblem{Class~A: median $82$, IQR $6$. Class~B: median $78$, IQR $14$. Which class performed better overall? Support your answer with data.}{class and reasoning}
\answerExplain{Class A}{Class A has a higher median ($82$ vs.\ $78$) and a smaller IQR ($6$ vs.\ $14$), meaning most students scored close to $82$. Class B is both lower and more spread out.}

\wordProblem{Group~1: mean $70$, median $72$, range $20$, IQR $8$. Group~2: mean $75$, median $74$, range $40$, IQR $18$. Write a comparison using at least two measures.}{comparison}
\answerExplain{Group 2 scores higher on average but is less consistent}{Group 2's median ($74$) and mean ($75$) are both higher. However, Group 1 is more consistent: IQR of $8$ vs.\ $18$ and range of $20$ vs.\ $40$.}

\wordProblem{Two basketball teams record points per game over $10$ games. Team~A: mean $78$, range $12$, IQR $5$. Team~B: mean $80$, range $35$, IQR $16$. A coach wants the most \emph{reliable} team. Which should she pick?}{team and reasoning}
\answerExplain{Team A}{Although Team B scores slightly more on average, Team A is far more predictable: its IQR ($5$) and range ($12$) are much smaller. Reliability matters when planning strategy.}

\end{practiceBox}

% ── Challenge ────────────────────────────────────────────────────────────────
\begin{challengeBox}
\prob Create two data sets of $7$ values each that have the \emph{same} median but \emph{different} IQRs. Show your work. \answerBlank[5cm]
\answerExplain{Example: A $= \{18, 19, 20, 21, 22, 23, 24\}$, IQR $= 4$; B $= \{5, 10, 15, 21, 30, 35, 40\}$, IQR $= 25$}{Both have median $21$. A: $Q_1 = 19$, $Q_3 = 23$, IQR $= 4$. B: $Q_1 = 10$, $Q_3 = 35$, IQR $= 25$. Same center, very different spread.}

\prob A data set has mean $50$ and IQR $10$. You add the value $200$ to the set. Which measure changes more --- the mean or the IQR? Explain.
\answerExplain{The mean changes more}{Adding $200$ is an extreme outlier that pulls the mean upward significantly. IQR only depends on $Q_1$ and $Q_3$, which are resistant to outliers and barely change.}
\end{challengeBox}

\encouragement{You're a data analysis expert --- comparing populations with confidence using center and spread!}
